%% Section 2, A taxonomy of the problem. Converted from paper/sections/02_taxonomy.md.
\section{A taxonomy of the problem}
\label{sec:taxonomy}

\subsection{Task families}
\label{sec:taxonomy:tasks}

The six families below are read off the rows rather than imposed on them, and Sec.~2.1's closing
paragraph says what they miss. Grasping is the largest. Fifty-seven of the 112 method rows carry the
grasp label, the labels are not exclusive, and one paper can sit in several families. Success is a
lift that survives a hold, and the thresholds differ by more than an order of magnitude. DexGraspVLA
\cite{dexgraspvla_2025} requires the object ``held 10 cm above the table for 20 s'', while Omnigrasp
\cite{omnigrasp_2024} requires it ``held at least 0.5 s in simulation''. What makes grasping hard at
scale is the continuum of starting configurations. UniDexGrasp++ \cite{unidexgrasp_pp_2023} states
it in Sec.~4.3, that ``we are dealing with an infinite number of tasks considering the initial
object pose can change continuously''.

Functional and tool use is second with 45 rows. It is the family where a stable grasp can still be
the wrong answer. Dexterous Functional Grasping \cite{dexterous_functional_grasping_2023} gives the
case in Sec.~2.1, that ``grabbing a
hammer from the head or handle are both equally valid ways of using it'', and only one of the two
lets the tool be used. Success is defined against the tool's function, and the field has no shared
way to state it.

In-hand reorientation has 31 rows and the most settled criteria, because the community inherited one
number. OpenAI \cite{openai_dexterity_2018} declared the goal achieved below 0.4 rad of orientation
error, DeXtreme \cite{dextreme_2022} keeps 0.4 rad at test time against a 0.1 rad training
tolerance, and Eureka \cite{eureka_2023}
counts consecutive successes at 0.1 rad. A shared tolerance is not a shared protocol, because the
stopping rule differs. Visual Dexterity \cite{visual_dexterity_2022} measures error ``when the
controller predicts it has
reached the goal and stops'', which makes the policy a judge of its own trial.

Tracking a human reference has 12 rows and inverts the problem. The target is a trajectory rather
than an endpoint, so success is a per-frame error band. ManipTrans \cite{maniptrans_2025} requires
all four of 30
degrees of rotation, 3 cm of translation, 8 cm of joint error and 6 cm of fingertip error. The
difficulty is that the reference came from a human hand and was never dynamically feasible for the
robot. DexMachina \cite{dexmachina_2025} warns in App.~B.4 that scoring by timesteps inside the
thresholds makes
``the results highly sensitive to the threshold values''.

Bimanual coordination has 43 rows and handover has 10. Handover is the smallest family and the one
whose failure has a single moment, because the giver must release only after the receiver has the
object. Dynamic Handover \cite{dynamic_handover_2023} reports a hit rate above its success rate and
blames the gap in Sec.~5.4 on ``occasional challenges encountered during the grasping phase of the
catcher''. Not every paper here learns both sides. In \cite{dexterous_handover_2025} only the
receiver is learned and the giver is a scripted arm.

Those six families do not cover the corpus. Twenty-four of the 112 method rows carry a label from
outside them and eight carry no label from the six at all. Twenty are labelled \texttt{other},
mostly generalist policies evaluated on a task suite rather than on a dexterous task family, among
them \cite{pi0_2024,pi05_2025,pistar06_2025,openvla_2024} and \cite{gemini_robotics_15_2025}. Three
adjacent families are named here rather than absorbed. Locomanipulation is
\cite{humanplus_2024,omnih2o_2024} and \cite{groot_n16_2025}, where the base is not fixed and the
gravity argument below changes character. Piano playing is \cite{robopianist_2023,rp1m_2024} and
\cite{pianomime_2024}, discussed in Section 6, where success is a per-timestep F1 against a MIDI
score rather than an object pose. Deformable manipulation is \cite{dexdeform_2023}, whose object
carries its own state and its own physics. Ferrari and Canny's grasp-quality measure
\cite{ferrari_canny_1992} carries no task label, because it is a measure rather than a task.
Table~\ref{tab:tasks}'s six rows are the families with enough papers to compare, and not a partition
of the corpus.

\subsection{What makes dexterous control hard}
\label{sec:taxonomy:hard}

Contact is non-smooth, and that is a property of the problem rather than of any solver. Bicchi
\cite{bicchi_grasping_chapter_2001} states in Sec.~1.2 that contact constraints are unilateral, and
that losing a contact ``involves an abrupt change of the structure of the model under
consideration''. Its Sec.~1.4 gives the sharper version, that a rod sliding on rough ground has
configurations with no consistent solution and configurations with more than one. Simulators inherit
the difficulty. The re-implementation study \cite{contact_models_comparison_2023} concludes in
Sec.~V that ``there is no fully satisfactory approach at the moment, as all existing solutions
compromise either accuracy, robustness, or efficiency''. Pang et al.
\cite{pang_global_planning_2022} treat non-smoothness as the thing to remove, curating every contact
pair in Sec.~III-D so that contact points and normals change smoothly.

The hand has more actuators than the object has degrees of freedom and is still short of authority
over it. The object moves only through contacts, and the contacts are what cannot be commanded.
Bicchi \cite{bicchi_grasping_chapter_2001} states in Sec.~1.4 that in multifingered grasps ``the
number of independent contact forces is much larger than the number of actuators. Thus, from a
controllability standpoint, not all the contact forces are controllable''.

Vision is occluded by the hand doing the work. Bai et al. \cite{bai_unified_manip_survey_2025} name
it in Sec.~4.3,
that ``occlusion hampers object tracking''. DexPoint \cite{dexpoint_2022} keeps vision and adds
imagined hand points
to the point cloud. Yin et al. \cite{rotating_without_seeing_2023} remove vision entirely and rotate
objects from
16 binary touch sensors over the palm, links and fingertips.

Gravity direction changes the task rather than scaling it. Visual Dexterity
\cite{visual_dexterity_2022} states that reorientation with the hand below the object ``is much
easier'', because ``with a downward-facing hand, the hand must manipulate the object while
simultaneously counteracting gravity''. AnyRotate \cite{anyrotate_2024} makes the direction a
randomised variable by ``randomly initializing hand orientations between episodes'', and its
Sec.~5.3 measures the cost, with performance dropping progressively from palm up and palm down,
through base up and base down, to thumb up and thumb down.

\subsection{Single hand versus two}
\label{sec:taxonomy:twohands}

Three counts describe two hands and they measure different things. Fifty-three of the 112 method
rows record two hands on the robot, which is the \texttt{bimanual} field. Forty-three carry the
\texttt{bimanual-coord} task label, the narrower claim that coordinating the hands is the task. Section 6
narrows again, to the 28 rows whose notes place a learned closed-loop controller on two
multi-fingered hands, and its opening paragraph names every exclusion that takes the 53 down to
the 28. That 28 is the denominator for every architecture count in this survey. Every bimanual
claim here names which of the three it uses.

Four things genuinely change when the second hand arrives. Contact stays non-smooth, occlusion
stays, and gravity stays the same problem.

Role asymmetry is the first. AsymDex \cite{asymdex_2024} assigns a dominant hand with full finger
and wrist control and a facilitating hand with 6-DoF base pose only, so that ``the facilitating hand
repositions and reorients one object, while the dominant hand performs complex manipulations''.

The second is that two hands on one object close a kinematic chain through the object. In
\cite{bicchi_grasping_chapter_2001} Sec.~1.4 the hand and object dynamics are separate and are
``linked through the n rigid-body contact constraints''. A second hand adds a second constraint set
on the same object rather than a second independent problem. That is where physical plausibility
fails first. BimanGrasp \cite{bimangrasp_2024} is the corpus row that measures it, rejecting a
synthesised grasp when ``total penetrations exceeded 1.5 mm'' and reporting in Sec.~IV-C that
``penetration remains the primary cause of grasp failure''.

The third is two-arm collision, which does not exist for one hand. DyDexHandover
\cite{dydexhandover_2025} resets an episode when ``any unintended arm contact with the environment
or self-collisions'' occurs. Bunny-VisionPro \cite{bunny_visionpro_2024} pays for it in the
controller, where adding a sphere-approximated self-collision cost raises motion-control time from
0.74 ms to 7.85 ms. Bidex \cite{bidex_teleop_2024} solves it in hardware, mounting the hands so the
human arm and the teacher arm ``are perpendicular to each other and do not collide''. DexMimicGen
\cite{dexmimicgen_2024} states plainly that it does not handle inter-arm collision at all.

The fourth is the doubled action space. Bi-DexHands \cite{bidexhands_2022} runs 20 tasks on two
Shadow Hands and finds single-agent PPO beating multi-agent RL on most of them, offering in Sec.~5.2
that ``PPO algorithm is able to use all observations for training the policy, while MARL can only
use partial observations''. AsymDex \cite{asymdex_2024} attacks the dimensionality from the other
side, halving the observation and action dimension through its role split and a frame relative to
the facilitating hand's object.

% ---------------------------------------------------------------------------
% Table 1 is hand-written rather than generated from corpus/rows, so it is
% carried in this section file instead of tex/tables/. Full text width.
% ---------------------------------------------------------------------------
\begin{table*}[!t]
  \caption{Task families against the properties that define success.}
  \label{tab:tasks}
  \centering\scriptsize
  \begin{tabular}{@{}L{0.085\textwidth}L{0.150\textwidth}L{0.150\textwidth}L{0.145\textwidth}L{0.130\textwidth}L{0.240\textwidth}@{}}
    \toprule
    \rowcolor{wash} \hdr{task family} & \hdr{what is held} & \hdr{what moves} & \hdr{what gravity does} & \hdr{what counts as failure} & \hdr{shortest honest success criterion} \\
    \midrule
    In-hand reorientation &
    one rigid object, held by the fingers for the whole episode &
    the object's orientation in the palm frame, while contacts break and remake &
    a disturbance whose direction is set by hand orientation, and a load to counteract with the palm down, \cite{visual_dexterity_2022} &
    the object leaves the hand, e.g.\ the fall reset at goal distance 0.24 m in \cite{dextreme_2022} &
    orientation error below a stated tolerance under a stated stopping rule, 0.4 rad in \cite{openai_dexterity_2018} and \cite{dextreme_2022}, 0.1 rad in \cite{eureka_2023} \\
    \addlinespace[2.5pt]
    Grasping &
    nothing at the start, the object rests on a support &
    the hand closes, then the object is lifted clear &
    the test itself, since the grasp exists to resist it &
    no lift, or a drop inside the hold window &
    a stated lift height held for a stated duration, 10 cm for 20 s in \cite{dexgraspvla_2025}, 10 cm to episode end in \cite{graspxl_2024}, 0.5 s in \cite{omnigrasp_2024} \\
    \addlinespace[2.5pt]
    Functional and tool use &
    a tool, held in the pose its function requires rather than the most stable one &
    the tool's own joint, or the tool acting on a second object &
    secondary, the function dominates &
    a stable grasp in the wrong place, \cite{dexterous_functional_grasping_2023} &
    no agreed criterion. Criteria range from never dropping the object, \cite{dexterous_functional_grasping_2023}, to a four-point rubric, \cite{dexvla_2025}. The one mechanical threshold is holding the tool joint past a commanded angle for the episode, \cite{articulated_tools_inhand_2025} \\
    \addlinespace[2.5pt]
    Tracking a human reference &
    whatever the human held in the captured clip &
    hand and object follow a recorded trajectory under physics &
    full, and the reference never obeyed it, so the reference is not dynamically feasible &
    per-frame error leaves the band, or the object drops &
    all stated per-frame errors below threshold for the whole clip, 30 deg and 3 cm and 8 cm joint and 6 cm fingertip in \cite{maniptrans_2025}. Thresholds are not comparable across papers, and DexMachina \cite{dexmachina_2025} calls the results ``highly sensitive to the threshold values'' \\
    \addlinespace[2.5pt]
    Bimanual coordination &
    one object by two hands, or one object per hand &
    both hands, the object, and the relative pose of the hands &
    can be carried by one hand while the other acts, \cite{asymdex_2024} &
    the object drops, the arms collide, or the hands work against each other &
    no agreed criterion. Task-specific thresholds dominate, e.g.\ 0.035 m block-to-cup distance in \cite{asymdex_2024}. The one criterion with a physical-plausibility gate is 2.0 s of hold with total penetration under 1.5 mm, \cite{bimangrasp_2024} \\
    \addlinespace[2.5pt]
    Handover and in-hand transfer &
    by the giver at the start, by the receiver at the end &
    the object crosses between hands, released only after it is regained &
    unopposed during the exchange, and total in a thrown transfer, \cite{dynamic_handover_2023} &
    a drop at the moment of release, and the catcher's grasp phase, \cite{dynamic_handover_2023} Sec.~5.4 &
    the receiver holds the object and moves it clear, more than 10 cm from the first hand in \cite{hato_visuotactile_2024}, or holds it to episode end in \cite{dydexhandover_2025} \\
    \bottomrule
  \end{tabular}
\end{table*}

Two cells say the field has no agreed criterion, and both are in families where a second body is
involved. A criterion naming a distance, a duration and a trial count can be re-run by someone else.
A rubric cannot, and Section 7 takes up what follows from that.
